% GENERATED FILE. Edit canonical lessons, then run npm run generate:books.
% canonical-source-hash: fnv1a64:848ffaf794192f04
% canonical-lessons: MR-C322-puraskar, MR-C322-satkar-karne, MR-C322-gaurav, MR-C322-abhinandan-karne, MR-C322-udghatan-karne

\chapter{Honours and Awards}
\label{ch:mr-honours-and-awards}
\hlchaptermodality{322}

\begin{chapteropening}
\textbf{By the end of this chapter:} \emph{I can talk about awards, honours and congratulations.}

Complete the last lesson of chapter 322: I can talk about awards, honours and congratulations.
\end{chapteropening}

\section[\texorpdfstring{\mr{पुरस्कार} (puraskār)}{पुरस्कार (puraskār)}]{\mr{पुरस्कार} (puraskār) — an award}

\label{lesson:MR-C322-puraskar}



\begin{quote}
Before the new one: say the Marathi for to attend, then the Marathi for attendance.
\end{quote}

\subsection*{You'll want to know: \mr{पुरस्कार}}
\textbf{\mr{पुरस्कार}} — \emph{puraskār} — \textquotedblleft{}an award\textquotedblright{}.

\begin{quote}
\mr{तिला} \mr{तिच्या} \mr{कामासाठी} \mr{मोठा} \mr{पुरस्कार} \mr{मिळाला}.
\end{quote}

\begin{quote}
\emph{tilā ticyā kāmāsāṭhī moṭhā puraskār miḷālā.}
\end{quote}

\begin{quote}
She received a big award for her work.
\end{quote}

Masculine.

\mr{पुरस्कार} \mr{मिळणे}, to receive an award, with the receiver in -\mr{ला}: \mr{तिला}. The letter \mr{ञ} hides inside \mr{ज्ञ}: \mr{ज्ञान}, knowledge.

\begin{cousinweb}
\textbf{\mr{पुरस्कार}} (\emph{puraskār}) is a tatsama, taken from Sanskrit \emph{puraskāra}, 'placing in front; honouring' (Monier-Williams).
\end{cousinweb}

\begin{grammarlens}[title={What you've built}]
One more word for telling what happened.
\end{grammarlens}

\subsection*{Your turn}
\begin{itemize}
\raggedright
  \item \emph{Say it:} \emph{puraskār}
  \item \emph{Say it:} \emph{puraskār}, once more
  \item \emph{Say it:} \emph{tilā ticyā kāmāsāṭhī moṭhā puraskār miḷālā.} — She received a big award for her work.
  \item \emph{Recall:} say the Marathi for to attend, then the Marathi for attendance, then say \emph{puraskār} again
\end{itemize}

\begin{tcolorbox}[breakable,colback=teal!4,colframe=teal!35!black,title={Before you move on}]
What does \mr{पुरस्कार} mean? (\textquotedblleft{}An award\textquotedblright{}.) Say it once more.
\end{tcolorbox}
\section[\texorpdfstring{\mr{सत्कार} \mr{करणे} (satkār karṇe)}{सत्कार करणे (satkār karṇe)}]{\mr{सत्कार} \mr{करणे} (satkār karṇe) — to felicitate; to honour publicly}

\label{lesson:MR-C322-satkar-karne}



\begin{quote}
Before the new one: say the Marathi for attendance, then the Marathi for an award.
\end{quote}

\subsection*{You'll want to know: \mr{सत्कार} \mr{करणे}}
\textbf{\mr{सत्कार} \mr{करणे}} — \emph{satkār karṇe} — \textquotedblleft{}to felicitate; to honour publicly\textquotedblright{}.

\begin{quote}
\mr{गावाने} \mr{विजेत्या} \mr{खेळाडूचा} \mr{सत्कार} \mr{केला}.
\end{quote}

\begin{quote}
\emph{gāvāne vijetyā kheḷāḍūcā satkār kelā.}
\end{quote}

\begin{quote}
The village honoured the winning player.
\end{quote}

Past, with -\mr{ने} (he …): \mr{त्याने} \mr{सत्कार} \mr{केला}, \emph{tyāne satkār kelā}.

The person honoured takes -\mr{चा}: \mr{शिक्षकांचा} \mr{सत्कार} \mr{करणे}; \mr{सत्कार} is masculine.

\begin{cousinweb}
\textbf{\mr{सत्कार} \mr{करणे}} (\emph{satkār karṇe}) joins \emph{satkār}, a tatsama taken from Sanskrit \emph{satkāra}, 'kind treatment, honour' (Monier-Williams), to \emph{karṇe}, 'to do'.
\end{cousinweb}

\begin{grammarlens}[title={What you've built}]
One more word for telling what happened.
\end{grammarlens}

\subsection*{Your turn}
\begin{itemize}
\raggedright
  \item \emph{Say it:} \emph{satkār karṇe}
  \item \emph{Say it:} \emph{satkār karṇe}, once more
  \item \emph{Say it:} \emph{gāvāne vijetyā kheḷāḍūcā satkār kelā.} — The village honoured the winning player.
  \item \emph{Recall:} say the Marathi for attendance, then the Marathi for an award, then say \emph{satkār karṇe} again
\end{itemize}

\begin{tcolorbox}[breakable,colback=teal!4,colframe=teal!35!black,title={Before you move on}]
What does \mr{सत्कार} \mr{करणे} mean? (\textquotedblleft{}To felicitate; to honour publicly\textquotedblright{}.) Say it once more.
\end{tcolorbox}
\section[\texorpdfstring{\mr{गौरव} (gaurav)}{गौरव (gaurav)}]{\mr{गौरव} (gaurav) — an honour; glory}

\label{lesson:MR-C322-gaurav}



\begin{quote}
Before the new one: say the Marathi for an award, then the Marathi for to felicitate.
\end{quote}

\subsection*{You'll want to know: \mr{गौरव}}
\textbf{\mr{गौरव}} — \emph{gaurav} — \textquotedblleft{}an honour; glory\textquotedblright{}.

\begin{quote}
\mr{हा} \mr{पुरस्कार} \mr{आमच्या} \mr{शाळेचा} \mr{गौरव} \mr{आहे}.
\end{quote}

\begin{quote}
\emph{hā puraskār āmacyā śāḷecā gaurav āhe.}
\end{quote}

\begin{quote}
This award is an honour for our school.
\end{quote}

Masculine.

\mr{गौरव} \mr{करणे}, to honour someone.

\begin{cousinweb}
\textbf{\mr{गौरव}} (\emph{gaurav}) is a tatsama, taken from Sanskrit \emph{gaurava}, 'weight; dignity, respect' (Monier-Williams).
\end{cousinweb}

\begin{grammarlens}[title={What you've built}]
One more word for telling what happened.
\end{grammarlens}

\subsection*{Your turn}
\begin{itemize}
\raggedright
  \item \emph{Say it:} \emph{gaurav}
  \item \emph{Say it:} \emph{gaurav}, once more
  \item \emph{Say it:} \emph{hā puraskār āmacyā śāḷecā gaurav āhe.} — This award is an honour for our school.
  \item \emph{Recall:} say the Marathi for an award, then the Marathi for to felicitate, then say \emph{gaurav} again
\end{itemize}

\begin{tcolorbox}[breakable,colback=teal!4,colframe=teal!35!black,title={Before you move on}]
What does \mr{गौरव} mean? (\textquotedblleft{}An honour; glory\textquotedblright{}.) Say it once more.
\end{tcolorbox}
\section[\texorpdfstring{\mr{अभिनंदन} \mr{करणे} (abhinandan karṇe)}{अभिनंदन करणे (abhinandan karṇe)}]{\mr{अभिनंदन} \mr{करणे} (abhinandan karṇe) — to congratulate}

\label{lesson:MR-C322-abhinandan-karne}



\begin{quote}
Before the new one: say the Marathi for to felicitate, then the Marathi for an honour.
\end{quote}

\subsection*{You'll want to know: \mr{अभिनंदन} \mr{करणे}}
\textbf{\mr{अभिनंदन} \mr{करणे}} — \emph{abhinandan karṇe} — \textquotedblleft{}to congratulate\textquotedblright{}.

\begin{quote}
\mr{सगळ्यांनी} \mr{पुरस्कार} \mr{मिळाल्याबद्दल} \mr{तिचं} \mr{अभिनंदन} \mr{केलं}.
\end{quote}

\begin{quote}
\emph{sagaḷyānnī puraskār miḷālyābaddal ticaṁ abhinandan kelaṁ.}
\end{quote}

\begin{quote}
Everyone congratulated her on receiving the award.
\end{quote}

Past, with -\mr{ने} (he …): \mr{त्याने} \mr{अभिनंदन} \mr{केलं}, \emph{tyāne abhinandan kelaṁ}.

The person takes -\mr{चं}: \mr{तिचं} \mr{अभिनंदन} \mr{केलं}; the reason takes -\mr{बद्दल}.

\begin{cousinweb}
\textbf{\mr{अभिनंदन} \mr{करणे}} (\emph{abhinandan karṇe}) joins \emph{abhinandan}, a tatsama taken from Sanskrit \emph{abhinandana}, 'rejoicing at, praising' (Monier-Williams), to \emph{karṇe}, 'to do'.
\end{cousinweb}

\begin{grammarlens}[title={What you've built}]
One more word for telling what happened.
\end{grammarlens}

\subsection*{Your turn}
\begin{itemize}
\raggedright
  \item \emph{Say it:} \emph{abhinandan karṇe}
  \item \emph{Say it:} \emph{abhinandan karṇe}, once more
  \item \emph{Say it:} \emph{sagaḷyānnī puraskār miḷālyābaddal ticaṁ abhinandan kelaṁ.} — Everyone congratulated her on receiving the award.
  \item \emph{Recall:} say the Marathi for to felicitate, then the Marathi for an honour, then say \emph{abhinandan karṇe} again
\end{itemize}

\begin{tcolorbox}[breakable,colback=teal!4,colframe=teal!35!black,title={Before you move on}]
What does \mr{अभिनंदन} \mr{करणे} mean? (\textquotedblleft{}To congratulate\textquotedblright{}.) Say it once more.
\end{tcolorbox}
\section[\texorpdfstring{\mr{उद्घाटन} \mr{करणे} (udghāṭan karṇe)}{उद्घाटन करणे (udghāṭan karṇe)}]{\mr{उद्घाटन} \mr{करणे} (udghāṭan karṇe) — to inaugurate}

\label{lesson:MR-C322-udghatan-karne}



\begin{quote}
Before the new one: say the Marathi for an honour, then the Marathi for to congratulate.
\end{quote}

\subsection*{You'll want to know: \mr{उद्घाटन} \mr{करणे}}
\textbf{\mr{उद्घाटन} \mr{करणे}} — \emph{udghāṭan karṇe} — \textquotedblleft{}to inaugurate\textquotedblright{}.

\begin{quote}
\mr{महापौरांनी} \mr{नव्या} \mr{पुलाचं} \mr{उद्घाटन} \mr{केलं}.
\end{quote}

\begin{quote}
\emph{mahāpaurānnī navyā pulācaṁ udghāṭan kelaṁ.}
\end{quote}

\begin{quote}
The mayor inaugurated the new bridge.
\end{quote}

Past, with -\mr{ने} (he …): \mr{त्याने} \mr{उद्घाटन} \mr{केलं}, \emph{tyāne udghāṭan kelaṁ}.

What is opened takes -\mr{चं}: \mr{शाळेचं} \mr{उद्घाटन} \mr{करणे}.

\begin{cousinweb}
\textbf{\mr{उद्घाटन} \mr{करणे}} (\emph{udghāṭan karṇe}) joins \emph{udghāṭan}, a tatsama taken from Sanskrit \emph{udghāṭana}, 'opening, unlocking' (Monier-Williams), to \emph{karṇe}, 'to do'.
\end{cousinweb}

\begin{grammarlens}[title={What you've built}]
One more word for telling what happened.
\end{grammarlens}

\subsection*{Your turn}
\begin{itemize}
\raggedright
  \item \emph{Say it:} \emph{udghāṭan karṇe}
  \item \emph{Say it:} \emph{udghāṭan karṇe}, once more
  \item \emph{Say it:} \emph{mahāpaurānnī navyā pulācaṁ udghāṭan kelaṁ.} — The mayor inaugurated the new bridge.
  \item \emph{Recall:} say the Marathi for an honour, then the Marathi for to congratulate, then say \emph{udghāṭan karṇe} again
\end{itemize}

\begin{tcolorbox}[breakable,colback=teal!4,colframe=teal!35!black,title={Before you move on}]
What does \mr{उद्घाटन} \mr{करणे} mean? (\textquotedblleft{}To inaugurate\textquotedblright{}.) Say it once more.
\end{tcolorbox}
